% Part: applied-modal-logics
% Chapter: epistemic-logic
% Section: language-epistemic-logic

\documentclass[../../../include/open-logic-section]{subfiles}

\begin{document}

\olfileid[zh]{aml}{el}{pal}

\olsection{公开宣告逻辑}

动态认知逻辑使我们得以刻画主体知识随时间流逝、或随着新知获得而变化的方式。其中许多逻辑借助信息性的\emph{事件}或\emph{更新}来刻画这些变化。最基本的更新类型是公开宣告，即某个公式被如实地宣告，且所有主体共同见证这一过程。为此，我们如下地扩张语言：

\begin{defn}
设 $G$ 是主体符号的集合。带公开宣告的多主体认知逻辑的基本语言包含
\begin{enumerate}
  \tagitem{prvFalse}{!!{falsity}的命题常元~$\lfalse$。}{}
  \tagitem{prvTrue}{!!{truth}的命题常元~$\ltrue$。}{}
  \item 一个!!{denumerable}!!{propositional variable}集合：$\Obj
    p_0$, $\Obj p_1$, $\Obj p_2$, \dots
  \item 命题联结词：\startycommalist
  \iftag{prvNot}{\ycomma $\lnot$（否定）}{}%
  \iftag{prvAnd}{\ycomma $\land$（合取）}{}%
  \iftag{prvOr}{\ycomma $\lor$（析取）}{}%
  \iftag{prvIf}{\ycomma $\lif$（!!{conditional}）}{}%
  \item 知识算子 $\Knows_a$，其中 $a \in G$。
  \item 公开宣告算子 $[!B]$，其中 $!B$ 是!!a{formula}。
\end{enumerate}
\end{defn}

公开宣告算子的作用与盒子算子相似，因此我们对语言的归纳定义也随之给出：

\begin{defn}
  认知语言的\emph{!!^{formula}}按如下方式归纳地定义：
  \begin{enumerate}
  \tagitem{prvFalse}{$\lfalse$ 是原子!!{formula}。}{}

  \tagitem{prvTrue}{$\ltrue$ 是原子!!{formula}。}{}

  \item 每个命题变元 $\Obj p_i$ 都是（原子的）
  !!{formula}。

  \tagitem{prvNot}{如果 $!A$ 是!!a{formula}，则 $\lnot !A$ 是
  !!a{formula}。}{}

  \tagitem{prvAnd}{如果 $!A$ 和 $!B$ 都是!!{formula}，则 $(!A \land
  !B)$ 是!!a{formula}。}{}

  \tagitem{prvOr}{如果 $!A$ 和 $!B$ 都是!!{formula}，则 $(!A \lor !B)$
  是!!a{formula}。}{}

  \tagitem{prvIf}{如果 $!A$ 和 $!B$ 都是!!{formula}，则 $(!A \lif !B)$
  是!!a{formula}。}{}

  \tagitem{prvIff}{如果 $!A$ 和 $!B$ 都是!!{formula}，则 $(!A \liff !B)$
  是!!a{formula}。}{}

  \item 如果 $!A$ 是!!a{formula}且 $a \in G$，则 $\Knows_a !A$ 是
  !!a{formula}。

  \item 如果 $!A$ 和 $!B$ 都是!!{formula}，则 $[!A] !B$ 是!!a{formula}。

  \tagitem{limitClause}{除此之外没有其他!!a{formula}。}{}
\end{enumerate}
\end{defn}

!!{formula} $[!A] !B$ 的意思是「在 $!A$ 被如实宣告之后，$!B$~成立」。有时在公开宣告的语境下讨论公共知识也会很有用，因此语言还可以包含公共知识算子~$\CKnows_G !A$。

\end{document}
